\documentclass[11pt,a4paper,twocolumn]{article}
\usepackage[T1]{fontenc}
\usepackage[utf8]{inputenc}
\usepackage[brazil]{babel}
\usepackage{lipsum}
\usepackage{hyperref}
\title{Prova do artigo sobre exportacao}
\author{Maria da Silva}
\date{11 de novembro de 2026}
\begin{document}
\maketitle
\section{Introducao}
Este artigo prova a exportacao para Word com duas colunas. \lipsum[1-2]
\section{Metodo}\label{sec:metodo}
O metodo usa uma nota de rodape.\footnote{Texto da nota de rodape do artigo.} \lipsum[3]
\begin{itemize}
\item primeiro item da lista
\item segundo item da lista
\item terceiro item da lista
\end{itemize}
\lipsum[4-5]
\subsection{Detalhes}
Veja o endereco \url{https://exemplo.org/artigo} para mais dados. \lipsum[6-7]
\section{Conclusao}
A conclusao fecha o texto. \lipsum[8-9]
\end{document}
